% ==================================================================
% id:       seq-011
% topic:    סדרות
% subtopic: סדרות מונוטוניות וסדרות נסיגה
% solution: full
% ==================================================================
\textbf{הוכיחו כי הסדרות הבאות מתכנסות וחשבו את גבולותיהן:}
\bE
\item
  $.a_1 = \sqrt2, \qquad a_{n+1} = \sqrt{2a_n}$
  
\hint{הוכיחו באינדוקציה מונוטוניות וחסימות. כדי לנחש חסם, הניחו שהגבול $L$ קיים והציבו אותו בשני אגפי נוסחת הנסיגה.}
\begin{solbox}
  ניתן להסתכל על מספר האיברים הראשונים בסדרה, ומכך לחשוד שהיא מונוטונית עולה וחסומה מלעיל על ידי
  $2$.
נוכיח באינדוקציה כי הסדרה הנתונה עולה וחסומה ע"י
$2$.

נחשוף בסוף הפתרון מדוע נלקח דווקא החסם הזה.
נתחיל בהוכחת החסימות.
\begin{itemize}
    \item שלב הבסיס:
מתקיים
$a_1=\sqrt{2}\leq2$.
\item שלב ההנחה:
נניח כי
$a_n\leq 2$
לאיזשהו
$n\geq 1$.
\item שלב ההיסק:
מתקיים
$$.a_{n+1} = \sqrt{2a_{n}}
\mathrel{\mathop{\leq}\limits_{\makebox[0pt][c]{\scriptsize\text{האינדוקציה הנחת}}}}\sqrt{2\cdot2} = 2$$
\end{itemize}
כך השלמנו את הוכחת החסימות.
נעבור עתה להוכחת המונוטוניות.
\begin{itemize}
    \item שלב הבסיס:
$$,a_2 = \sqrt{2\sqrt{2}} \geq  \sqrt{2}=a_1$$
כך וידאנו את שלב הבסיס של הוכחת המונוטוניות.
\item שלב ההנחה:
נניח כי
$$a_{n+1} \geq a_n$$
לאיזשהו
$n\geq 1$.
\item שלב ההיסק:
מתקיים
$$.a_{n+2} = \sqrt{2a_{n+1}}\mathrel{\mathop{\geq}\limits_{\makebox[0pt][c]{\scriptsize\text{האינדוקציה הנחת}}}}\sqrt{2a_n} = a_{n+1}$$
\end{itemize}
כך הושלמה ההוכחה של המונוטוניות.
נסיק כי הסדרה הנתונה מתכנסת
ונסמן את גבולה ב-
$L$.
מכיוון ש-
$$,a_{n+1}=\sqrt{2a_n}$$
מיחידות הגבול
מתקיים
$$,L=\lim_{n\to\infty}a_{n+1}=\lim_{n\to\infty}\sqrt{2a_n}=\sqrt{2L}$$
כלומר
$L^2=2L$.
ישנם שני פתרונות למשוואה הזו -
$L_1=0$
ו-
$L_2=2$.
מכיוון שהסדרה הנתונה חיובית ועולה, לא יתכן שגבולה הוא
$0$,
לכן
$L=2$.

$\bullet$
שימו לב שאת כל החלק האחרון -- מהשלב בו ידועה ההתכנסות ומחפשים את הגבול -- היה ניתן לעשות גם בתחילת הפתרון, כדי לקבל את החשוד או חשודים בתור הגבול של הסדרה, ולנסות להיעזר בחשוד זה כחסם עבור הסדרה, כמו שכאן הגבול הוא
$2$
והיא מונוטונית עולה אז
$2$
משמש גם כחסם מלעיל שלה.
הערה: במעבר לגבול השתמשנו באריתמטיקה של גבולות (גבול של סכום ושל מכפלה בקבוע), ובאריתמטיקה של חזקות: אם $b_n\goes b$ ו-$b_n,b\ge0$, אז $\sqrt{b_n}=b_n^{1/2}\goes b^{1/2}=\sqrt b$. כך אפשר להעביר את הגבול אל תוך השורש.
\end{solbox}

\item
$
  .a_1 = \sqrt{7}, \qquad a_{n+1} = \sqrt{7+a_n}
$
\hint{לחסימות, נסו להוכיח באינדוקציה חסם כללי $M$ ובדקו איזה תנאי על $M$ נדרש כדי שצעד האינדוקציה יעבוד.}
\begin{solbox}
נוכיח באינדוקציה כי הסדרה עולה וחסומה.
נתחיל בהוכחת המונוטוניות באופן מאוד דומה לסעיף הקודם.
\begin{itemize}
    \item שלב הבסיס:
    מתקיים
    $a_1=\sqrt{7}<\sqrt{7+\sqrt{7}}=a_2$.
    \item שלב ההנחה:
    נניח כי
    $a_n<a_{n+1}$
    לאזשהו
    $n\geq 1$.
    \item
    שלב ההיסק:
    מתקיים
    $$.a_{n+1}=\sqrt{7+a_n}\mathrel{\mathop{<}\limits_{\makebox[0pt][c]{\scriptsize\text{היצקודניאה תחנה}}}}\sqrt{7+a_{n+1}}=a_{n+2}$$
\end{itemize}
    בכך הושלמה הוכחת המונוטוניות. נעבור להוכחת החסימות. אך אנו לא יודעים מה יהיה חסם של הסדרה. בפתרון זה נציג כיצד ניתן להבין מה יהיה חסם נכון תוך כדי פתרון ללא ידיעתו מראש!

    נתחיל להוכיח באינדוקציה כי הסדרה חסומה ע"י
    $M$
    ותוך כדי ההוכחה נאסוף עוד תנאים על
    $M$,
    תנאים שיגרמו להוכחה להיות נכונה ושלא יכולנו לדעת אותם לפני כן.
\begin{itemize}
    \item שלב הבסיס:
    $,a_1=\sqrt{7}\leq M$
    מכך רואים כי הדרישה הראשונה על
    $M$
    היא  שיקיים
    $M\geq\sqrt{7}$.
    \item שלב ההנחה:
    נניח כי
    $a_n\leq M$
    לאזשהו
    $n\geq 1$.
    \item
    שלב ההיסק:
    מתקיים
    $$
a_{n+1} = \sqrt{7 + a_n}
\qquad\mathrel{\mathop{\leq}\limits_{\makebox[0pt][c]{\scriptsize\text{האינדוקציה הנחת}}}}\qquad
\sqrt{7 + M}
$$
כל מה שנרצה זה שיתקיים
$\sqrt{7+M}\leq M$
ואז הוכחת האינדוקציה תיגמר.
לשם כך נבין מה זה אומר ונוסיף זאת כאילוץ בבחירת
$M$.
מספיק למצוא
$M$
אחד המקיים זאת, אין צורך לפתור את האי-שיוויון באופן מלא. למשל,
$M=9$
מקיים
$\sqrt{7+M}=\sqrt{16}=4\le 9=M$
לכן זו תהיה בחירתנו, כי היא גם מקיימת את האילוץ הראשון שהיה
$M\ge \sqrt{7}$.
    כך השלמנו את הוכחת החסימות.

    נסיק כי הסדרה הנתונה מתכנסת ונסמן את גבולה ב-
    $L$.
    כמו בסעיף הקודם,
    מכיוון ש-
    $a_{n+1}=\sqrt{7+a_n}$
    נסיק מיחידות הגבול כי
    $$,L=\lim_{n\to\infty}a_{n+1}=\lim_{n\to\infty}\sqrt{7+a_n}=\sqrt{7+L}$$
    קיבלנו משוואה ריבועית
    $L^2-L-7=0$,
    לה שני פתרונות
    $L_1=\frac{1-\sqrt{29}}{2}$
    ו-
    $L_2=\frac{1+\sqrt{29}}{2}$.
    בבירור
    $L_1$
    לא מתאים כי
    $L_1<0$
    ואיברי הסדרה הם מספרים חיוביים.
    לכן
    $L=L_2$,
     כלומר
     $$.a_n\goes \frac{1+\sqrt{29}}{2}$$
     שימו לב שהפעם במהלך הפתרון השתמשנו בחסם
     $9$
     שהוא יותר גדול מהגבול של הסדרה, וזה בסדר גמור, כי רק רצינו להראות שהיא חסומה אז החסם לא חייב להיות החסם הכי טוב.
     \end{itemize}
הערה: במעבר לגבול השתמשנו באריתמטיקה של גבולות (גבול של סכום ושל מכפלה בקבוע), ובאריתמטיקה של חזקות: אם $b_n\goes b$ ו-$b_n,b\ge0$, אז $\sqrt{b_n}=b_n^{1/2}\goes b^{1/2}=\sqrt b$. כך אפשר להעביר את הגבול אל תוך השורש.
\end{solbox}

\item
$a_{n+1}=\frac{a_n+1}{2}$,
כאן
$a_1$
לא ידוע
(הכוונה: הפרידו למקרים עבור
$a_1$)
\hint{מצאו את המספר שעבורו $a_{n+1}=a_n$, ופצלו למקרים לפי מיקומו של $a_1$ ביחס אליו. בכל מקרה בדקו מונוטוניות וחסימות.}
\begin{solbox}
נפריד למקרים:
\begin{itemize}
\item
אם
$a_1=1$
אז ניתן לראות כי זוהי הסדרה הקבועה ולכן
מתכנסת ל-$1$.
\item
אם
$a_1>1$,
אז נראה כי זוהי סדרה מונוטונית יורדת החסומה מלרע על ידי
$0$.

ראשית, נוכיח באינדוקציה כי
$a_n>0$
לכל
$n\in\nn$:

$\bullet$
עבור
$n=1$
מתקיים
$a_1>1>0$;

$\bullet$
נניח כי
$a_n>0$
ולכן
$$a_{n+1}=
\frac{a_n+1}{2}>\frac{1}{2}>0$$
כפי שרצינו.

שנית, נוכיח באינדוקציה כי
$a_{n+1}<a_n$
לכל
$n\in\nn$:

$\bullet$
עבור
$n=1$
מתקיים
$a_2=\frac{a_1+1}{2}<a_1$
כי
$a_1>1$;

$\bullet$
נניח כי
$a_{n+1}<a_n$
ולכן
$$a_{n+2}=\frac{a_{n+1}+1}{2}<\frac{a_n+1}{2}=a_{n+1}$$
כפי שרצינו.
\item
אם
$a_1<1$,
אז נראה כי זוהי סדרה מונוטונית עולה החסומה מלעיל על ידי
$1$.

ראשית, נוכיח באינדוקציה כי
$a_n<1$
לכל
$n\in\nn$:

$\bullet$
עבור
$n=1$
מתקיים
$a_1<1$;

$\bullet$
נניח כי
$a_n<1$
ולכן
$a_{n+1}=\frac{a_n+1}{2}<1$,
כפי שרצינו.

שנית- והפעם ניתן לוותר על האינדוקציה ולהראות באופן ישיר- הסדרה מונוטונית עולה, כי לכל
$n\in\nn$
מתקיים
$$.a_{n+1}=\frac{a_n+1}{2}>\frac{a_n+a_n}{2}=a_n$$
\end{itemize}

בשני המקרים האחרונים קבלנו כי הסדרה חסומה ומונטונית, לכן היא מתכנסת, כעת נמצא את גבולה.
נסמן
$a_n\goes L$
ועל ידי שימוש באריתמטיקה נקבל כי
$$\frac{a_n+1}{2}\goes\frac{L+1}{2}$$
מצד שני, מתקיים
$a_{n+1}\goes L$
אבל
$a_{n+1}=\frac{a_n+1}{2}$
ולכן מיחידות הגבול (של סדרה מתכנסת) נקבל כי
$$L=\frac{L+1}{2}$$
כלומר $L=1$.
\end{solbox}
\eE
